% =============================================================================
% Appendix D — Color Table
% =============================================================================
\chapter{Color Table}

\epigraph{``Sixteen colours ought to be enough for anybody.''}%
{\textit{--- Every 8-bit computer designer, 1977--1985}}

\section*{Introduction}

The NovaVM uses a fixed 16-colour palette inspired by the Commodore~64.  Every
colour index in the system --- text foreground and background, graphics pixels,
and sprite pixels --- references this palette via a 4-bit value (0--15).

% ── Define all 16 colours ─────────────────────────────────────────────────────
\definecolor{nova0}{RGB}{0,0,0}
\definecolor{nova1}{RGB}{255,255,255}
\definecolor{nova2}{RGB}{136,0,0}
\definecolor{nova3}{RGB}{170,255,238}
\definecolor{nova4}{RGB}{204,68,204}
\definecolor{nova5}{RGB}{0,204,85}
\definecolor{nova6}{RGB}{0,0,170}
\definecolor{nova7}{RGB}{238,238,119}
\definecolor{nova8}{RGB}{221,136,85}
\definecolor{nova9}{RGB}{102,68,0}
\definecolor{nova10}{RGB}{255,119,119}
\definecolor{nova11}{RGB}{51,51,51}
\definecolor{nova12}{RGB}{119,119,119}
\definecolor{nova13}{RGB}{170,255,102}
\definecolor{nova14}{RGB}{0,136,255}
\definecolor{nova15}{RGB}{187,187,187}


% ═══════════════════════════════════════════════════════════════════════════════
\section{Palette Reference}
\label{sec:color-palette}
% ═══════════════════════════════════════════════════════════════════════════════

\begin{center}
\small
\begin{tabular}{@{}c c l r r r l@{}}
\toprule
\textbf{Index} & \textbf{Swatch} & \textbf{Name} &
\textbf{R} & \textbf{G} & \textbf{B} & \textbf{Hex} \\
\midrule
 0 & \tikz\draw[fill=nova0,  draw=retrogray] (0,0) rectangle (1.2em,1.2em); & Black       &   0 &   0 &   0 & \texttt{\#000000} \\
 1 & \tikz\draw[fill=nova1,  draw=retrogray] (0,0) rectangle (1.2em,1.2em); & White       & 255 & 255 & 255 & \texttt{\#FFFFFF} \\
 2 & \tikz\draw[fill=nova2,  draw=retrogray] (0,0) rectangle (1.2em,1.2em); & Red         & 136 &   0 &   0 & \texttt{\#880000} \\
 3 & \tikz\draw[fill=nova3,  draw=retrogray] (0,0) rectangle (1.2em,1.2em); & Cyan        & 170 & 255 & 238 & \texttt{\#AAFFEE} \\
 4 & \tikz\draw[fill=nova4,  draw=retrogray] (0,0) rectangle (1.2em,1.2em); & Purple      & 204 &  68 & 204 & \texttt{\#CC44CC} \\
 5 & \tikz\draw[fill=nova5,  draw=retrogray] (0,0) rectangle (1.2em,1.2em); & Green       &   0 & 204 &  85 & \texttt{\#00CC55} \\
 6 & \tikz\draw[fill=nova6,  draw=retrogray] (0,0) rectangle (1.2em,1.2em); & Blue        &   0 &   0 & 170 & \texttt{\#0000AA} \\
 7 & \tikz\draw[fill=nova7,  draw=retrogray] (0,0) rectangle (1.2em,1.2em); & Yellow      & 238 & 238 & 119 & \texttt{\#EEEE77} \\
 8 & \tikz\draw[fill=nova8,  draw=retrogray] (0,0) rectangle (1.2em,1.2em); & Orange      & 221 & 136 &  85 & \texttt{\#DD8855} \\
 9 & \tikz\draw[fill=nova9,  draw=retrogray] (0,0) rectangle (1.2em,1.2em); & Brown       & 102 &  68 &   0 & \texttt{\#664400} \\
10 & \tikz\draw[fill=nova10, draw=retrogray] (0,0) rectangle (1.2em,1.2em); & Light Red   & 255 & 119 & 119 & \texttt{\#FF7777} \\
11 & \tikz\draw[fill=nova11, draw=retrogray] (0,0) rectangle (1.2em,1.2em); & Dark Grey   &  51 &  51 &  51 & \texttt{\#333333} \\
12 & \tikz\draw[fill=nova12, draw=retrogray] (0,0) rectangle (1.2em,1.2em); & Medium Grey & 119 & 119 & 119 & \texttt{\#777777} \\
13 & \tikz\draw[fill=nova13, draw=retrogray] (0,0) rectangle (1.2em,1.2em); & Light Green & 170 & 255 & 102 & \texttt{\#AAFF66} \\
14 & \tikz\draw[fill=nova14, draw=retrogray] (0,0) rectangle (1.2em,1.2em); & Light Blue  &   0 & 136 & 255 & \texttt{\#0088FF} \\
15 & \tikz\draw[fill=nova15, draw=retrogray] (0,0) rectangle (1.2em,1.2em); & Light Grey  & 187 & 187 & 187 & \texttt{\#BBBBBB} \\
\bottomrule
\end{tabular}
\end{center}


% ═══════════════════════════════════════════════════════════════════════════════
\section{Usage Notes}
\label{sec:color-usage}
% ═══════════════════════════════════════════════════════════════════════════════

\begin{itemize}
  \item \textbf{Transparency:} Colour index~0 (black) serves as the transparent
        colour for sprites.  Sprite pixels with value~0 allow the background to
        show through.
  \item \textbf{Default colours:} At power-on, the text foreground is~1 (white)
        and the background is~0 (black).
  \item \textbf{Boot screen:} The EhBASIC cold-start routine sets the
        background to~6 (blue) and the foreground to~1 (white), mimicking the
        classic Commodore~64 appearance.
  \item \textbf{Colour RAM:} In text mode, each character cell has an associated
        byte in colour RAM.  The low nibble sets the foreground colour and the
        high nibble sets the background colour for that cell.
  \item \textbf{Graphics bitmap:} Each pixel byte uses only the low nibble
        (bits 3--0) for the colour index.  The upper nibble is ignored.
\end{itemize}


% ═══════════════════════════════════════════════════════════════════════════════
\section{BASIC Example}
\label{sec:color-example}
% ═══════════════════════════════════════════════════════════════════════════════

\begin{lstlisting}
10 REM cycle through all 16 colors
20 FOR C=0 TO 15
30   COLOR 1, C : CLS
40   PRINT "COLOR"; C
50   FOR D=1 TO 2000 : NEXT D
60 NEXT C
70 COLOR 1, 6 : CLS
\end{lstlisting}
